\answer
\begin{enumerate}

%% 4章 問1（4.2節）
\item[4.2節 問1]
H $=$ V$\cdot$s/Aだから
\begin{equation*}
[LI^2] = \frac{\mathrm{V\cdot s}}{\mathrm{A}}\times\mathrm{A}^2
= \mathrm{V\cdot A\cdot s} = \mathrm{W\cdot s} = \mathrm{J}
\end{equation*}
となり，確かにエネルギーの単位になる。
電圧×電流＝電力（W），電力×時間＝エネルギー（J）という筋道を単位がそのまま映している。

%% 4章 問2（4.2節）
\item[4.2節 問2]
式\ref{eq4.5}より，電圧一定なら電流の増加分は電圧×時間を$L$で割ったものである。
\begin{equation*}
\mathit{\Delta}i_L = \frac{V\,\mathit{\Delta}t}{L}
= \frac{5\times10\times10^{-6}}{100\times10^{-6}} = 0.5\,\mathrm{A}
\end{equation*}

%% 4章 問3（4.3節）
\item[4.3節 問3]
式\ref{eq4.9}より
\begin{equation*}
\mathit{\Delta}v_C = \frac{I\,\mathit{\Delta}t}{C}
= \frac{1\times10\times10^{-6}}{10\times10^{-6}} = 1\,\mathrm{V}
\end{equation*}
\ref{sec4.2}節の問2の$V \leftrightarrow I$，$L \leftrightarrow C$，
$\mathit{\Delta}i \leftrightarrow \mathit{\Delta}v$を入れ替えた双対の問題になっており，
計算の形はまったく同じである。

%% 4章 問4（4.4節）
\item[4.4節 問4]
$LC$の単位は
\begin{equation*}
[LC] = \frac{\mathrm{V\cdot s}}{\mathrm{A}}\times\frac{\mathrm{A\cdot s}}{\mathrm{V}}
= \mathrm{s}^2
\end{equation*}
だから$\sqrt{LC}$の単位はsであり，
$f_0 = 1/(2\pi\sqrt{LC})$の単位は$1/\mathrm{s} = \mathrm{Hz}$になる。

%% 4章 問5（4.5節）
\item[4.5節 問5]
$\omega$の単位は$1/\mathrm{s}$だから
\begin{equation*}
[\omega L] = \frac{1}{\mathrm{s}}\times\frac{\mathrm{V\cdot s}}{\mathrm{A}}
= \frac{\mathrm{V}}{\mathrm{A}} = \Omega，\qquad
\left[\frac{1}{\omega C}\right]
= \mathrm{s}\times\frac{\mathrm{V}}{\mathrm{A\cdot s}}
= \frac{\mathrm{V}}{\mathrm{A}} = \Omega
\end{equation*}
どちらも電圧／電流の単位であり，抵抗と同じ土俵で比べられる。

\end{enumerate}
